% Preserve original Unicode in PDF copy/search when OpenType features select
% CID glyphs (oldstyle figures, superscripts, small caps and ligatures). This
% native XeTeX output switch records ActualText without changing shaped text.
\XeTeXgenerateactualtext=1

% PTXprint selects an empty language; use XeTeX's preloaded US English patterns.
% This also sets the standard minimum fragments: 2 letters left, 3 right.
\uselanguage{USenglish}

% Give note origin numbers a compact fixed 1 pt separation in the margin.
% Remove the source's trailing space/NBSP inside the origin style, then add
% the kern after the style closes. Keeping both hooks avoids an extra word
% space or a gap that changes during justification; retain the no-break tie.
% Register hooks through PTXprint's API; defining end-/after- macros alone
% does not put them in the style engine's hook lookup.
\sethook{end}{fr}{\unskip}
\sethook{end}{xo}{\unskip}
\sethook{after}{fr}{\kern1pt\nobreak}
\sethook{after}{xo}{\kern1pt\nobreak}

% Put cross references and footnotes beside their source verses in the wide
% inner margin, using PTXprint's native side-aligned notes rather than a
% second scripture column. Its repeated runs adjust crowded note positions;
% the source reference and caller still identify their verse. This native
% marginal-note mode does not move overflow to the bottom of the page.
% Divide the 42 mm inner margin into a 6 mm gap,
% 24 mm of notes, and 12 mm clear at the binding. PTXprint's
% marginnoteswidth includes the gap, so the allocation is 30 mm, not 24 mm.
\newif\ifBibleMarginNote
\begingroup
\catcode`\@=11
\marginaln@te{x}
\marginaln@te{f}
\global\marginnotesgap=6mm
\global\marginnoteswidth=\dimexpr 24mm+\marginnotesgap\relax
\gdef\MarginNoteSide{inner}
\gdef\BibleStartNote#1{%
  \BibleMarginNotetrue
  \ifx\ch@pter\empty\BibleMarginNotefalse\fi
  \ifx\ch@pter\z@ro\BibleMarginNotefalse\fi
  \ifhe@dings\BibleMarginNotefalse\fi
  \ifBibleMarginNote
    \n@test@rt{#1}%
  \else
    % Separate paragraphs keep front-matter/heading notes readable;
    % PTXprint's normal insertions reserve space and can split long notes.
    \global\p@ranotesfalse
    \expandafter\insert\csname note-#1\endcsname
  \fi
}
\gdef\BibleEndNote#1{%
  \ifBibleMarginNote
    % Reuse PTXprint's globally incremented note ID: its marginal occurrence
    % counter is local and can repeat IDs for notes on the same verse.
    % n@te@nd increments this seed once; the resulting suffix stays unique.
    \expandafter\global\expandafter\let\csname nm-cref-#1\endcsname\c@rref
    \expandafter\global\csname nm-count-#1\endcsname=\n@teid\relax
    \n@te@nd{#1}%
  \fi
}
\expandafter\gdef\csname note-no-insert-f\endcsname{\BibleStartNote{f}}
\expandafter\gdef\csname note-no-insert-x\endcsname{\BibleStartNote{x}}
\expandafter\gdef\csname endnote-no-insert-f\endcsname{\BibleEndNote{f}}
\expandafter\gdef\csname endnote-no-insert-x\endcsname{\BibleEndNote{x}}
\endgroup

% Introduction notes use the native ef class in the generated USFM. Leaving
% that class out of marginaln@te keeps them below without a placement lookup.
% Including ef enables study-note formatting in the template; use its native
% switch to retain ordinary full-width bottom notes for these introductions.
% NotStudyNotes assigns locally, so call it outside the temporary group above.
\NotStudyNotes{ef}

% BSB numbers ordinary/no-verse pages in the center header and opening
% pages in the center footer. Put every visible folio at the outer bottom
% corner for a consistent locator on facing pages, clearing the old slots
% to avoid duplicates. Keep \pagenumber itself, so \nopagenums and
% \dopagenums govern
% suppression. PTXprint continues to choose page types and omit blank pages.
\def\RFoddright{\pagenumber}
\def\RFevenleft{\pagenumber}
\def\RFnoVoddright{\pagenumber}
\def\RFnoVevenleft{\pagenumber}
\let\RFtitleoddcenter\empty
\let\RFtitleevencenter\empty
\def\RFtitleoddright{\pagenumber}
\def\RFtitleevenleft{\pagenumber}

% Load the shared optical-margin adapter used by the font regression test.
\input ../../../shared/ptxprint/Bible/protrusion.tex
\input ../../../shared/ptxprint/Bible/word-spacing.tex

% Susanna and Bel and the Dragon have a zero-width-space \cp in place
% of a chapter
% number. Remove its drop-cap cutout so the opening text does not indent
% around a chapter figure that readers cannot see.
\begingroup
\catcode`\@=11
\xdef\BibleHiddenChapterText{\Uchar"200B}
\ifx\BibleMakeChapterBox\undefined\global\let\BibleMakeChapterBox\m@kechapterbox\fi
\gdef\m@kechapterbox{%
  \BibleMakeChapterBox
  \edef\BibleChapterText{\detokenize\expandafter{\ch@ptert@xt}}%
  \ifx\BibleChapterText\BibleHiddenChapterText
    \setbox\ch@pterbox=\hbox{}%
    \ch@pterwd=0pt
    \ch@pterdp=0pt
  \fi
}
\endgroup

% Share the reference-table style between the numbering and abbreviations pages.
% Select it at table entry, before PTXprint opens the table's group and reads
% the category. Row hooks run too late; transition hooks depend on the
% preceding paragraph marker. PTXprint clears the category after each table.
\newif\ifBibleReferenceTables
\setbookhook{start}{XXA}{\BibleReferenceTablestrue}
\setbookhook{end}{XXA}{\BibleReferenceTablesfalse}
\setbookhook{start}{XXD}{\BibleReferenceTablestrue}
\setbookhook{end}{XXD}{\BibleReferenceTablesfalse}
\begingroup
\catcode`\@=11
\global\let\BibleStartTable\startt@ble
\gdef\startt@ble{%
  \ifBibleReferenceTables
    \ifdoingt@ble\else\def\tablecategory{reference}\fi
  \fi
  \BibleStartTable
}
\endgroup

% Empty paragraphs do not enter the start hook. Put the rule in the
% heading box through the before hook, with fixed spacing in TeX points.
\sethook{before}{sd2}{%
  \vskip6pt
  \hbox to \hsize{\hfil\vrule width .25\hsize height .5pt depth 0pt\hfil}%
  \nobreak\vskip6pt
}
